% GENERATED FILE. Edit canonical lessons, then run npm run generate:books.
% canonical-source-hash: fnv1a64:96e24b90118bcc25
% canonical-lessons: HI-C435-rezgari, HI-C435-saiz, HI-C435-posak, HI-C435-rumal, HI-C435-simatna

\chapter{Change, Sizes and Clothes}
\label{ch:hi-change-sizes-and-clothes}
\hlchaptermodality{435}

\begin{chapteropening}
\textbf{By the end of this chapter:} \emph{I can ask for change, find my size, buy a handkerchief and say a child curled up against the cold.}

Complete the last lesson of chapter 435: I can ask for change, find my size, buy a handkerchief and say a child curled up against the cold.
\end{chapteropening}

\section[\texorpdfstring{\hi{रेज़गारी} (rezgārī)}{रेज़गारी (rezgārī)}]{\hi{रेज़गारी} (rezgārī) — loose coins, change}

\label{lesson:HI-C435-rezgari}



\begin{quote}
Before the new one: say the Hindi for to suit, then the Hindi for a touchstone.
\end{quote}

\subsection*{You'll want to know: \hi{रेज़गारी}}
\textbf{\hi{रेज़गारी}} — \emph{rezgārī} — \textquotedblleft{}loose coins, change\textquotedblright{}.

\begin{quote}
\hi{बस} \hi{के} \hi{लिए} \hi{रेज़गारी} \hi{रखिए।}
\end{quote}

\begin{quote}
\emph{bas ke liye rezgārī rakhiye.}
\end{quote}

\begin{quote}
Keep some loose change for the bus.
\end{quote}

Feminine.

\begin{cousinweb}
\textbf{\hi{रेज़गारी}} (\emph{rezgārī}) is usually traced to Persian \emph{reza}, \textquotedblleft{}a small piece, a fragment\textquotedblright{}: money broken into small pieces.
\end{cousinweb}

\begin{grammarlens}[title={What you've built}]
One more word for getting around.
\end{grammarlens}

\subsection*{Your turn}
\begin{itemize}
\raggedright
  \item \emph{Say it:} \emph{rezgārī}
  \item \emph{Say it:} \emph{rezgārī}, once more
  \item \emph{Say it:} \emph{bas ke liye rezgārī rakhiye.} — Keep some loose change for the bus.
  \item \emph{Recall:} say the Hindi for to suit, then the Hindi for a touchstone, then say \emph{rezgārī} again
\end{itemize}

\begin{tcolorbox}[breakable,colback=teal!4,colframe=teal!35!black,title={Before you move on}]
What does \hi{रेज़गारी} mean? (\textquotedblleft{}Loose coins, change\textquotedblright{}.) Say it once more.
\end{tcolorbox}
\section[\texorpdfstring{\hi{साइज़} (sāiz)}{साइज़ (sāiz)}]{\hi{साइज़} (sāiz) — a size}

\label{lesson:HI-C435-saiz}



\begin{quote}
Before the new one: say the Hindi for a touchstone, then the Hindi for loose coins.
\end{quote}

\subsection*{You'll want to know: \hi{साइज़}}
\textbf{\hi{साइज़}} — \emph{sāiz} — \textquotedblleft{}a size\textquotedblright{}.

\begin{quote}
\hi{इस} \hi{कमीज़} \hi{का} \hi{साइज़} \hi{क्या} \hi{है}?
\end{quote}

\begin{quote}
\emph{is kamīz kā sāiz kyā hai?}
\end{quote}

\begin{quote}
What size is this shirt?
\end{quote}

Masculine.

\begin{cousinweb}
\textbf{\hi{साइज़}} (\emph{sāiz}) is borrowed from English \emph{size}.
\end{cousinweb}

\begin{grammarlens}[title={What you've built}]
One more word for getting around.
\end{grammarlens}

\subsection*{Your turn}
\begin{itemize}
\raggedright
  \item \emph{Say it:} \emph{sāiz}
  \item \emph{Say it:} \emph{sāiz}, once more
  \item \emph{Say it:} \emph{is kamīz kā sāiz kyā hai?} — What size is this shirt?
  \item \emph{Recall:} say the Hindi for a touchstone, then the Hindi for loose coins, then say \emph{sāiz} again
\end{itemize}

\begin{tcolorbox}[breakable,colback=teal!4,colframe=teal!35!black,title={Before you move on}]
What does \hi{साइज़} mean? (\textquotedblleft{}A size\textquotedblright{}.) Say it once more.
\end{tcolorbox}
\section[\texorpdfstring{\hi{पोशाक} (pośāk)}{पोशाक (pośāk)}]{\hi{पोशाक} (pośāk) — a dress, an outfit}

\label{lesson:HI-C435-posak}



\begin{quote}
Before the new one: say the Hindi for loose coins, then the Hindi for a size.
\end{quote}

\subsection*{You'll want to know: \hi{पोशाक}}
\textbf{\hi{पोशाक}} — \emph{pośāk} — \textquotedblleft{}a dress, an outfit\textquotedblright{}.

\begin{quote}
\hi{मेले} \hi{में} \hi{बच्चे} \hi{नई} \hi{पोशाक} \hi{पहनकर} \hi{आए।}
\end{quote}

\begin{quote}
\emph{mele meṁ bacce naī pośāk pahankar āe.}
\end{quote}

\begin{quote}
The children came to the fair in new clothes.
\end{quote}

Feminine.

\begin{cousinweb}
\textbf{\hi{पोशाक}} (\emph{pośāk}) is usually traced to Persian \emph{pošāk}, \textquotedblleft{}clothing, dress\textquotedblright{}; it is feminine.
\end{cousinweb}

\begin{grammarlens}[title={What you've built}]
One more word for getting around.
\end{grammarlens}

\subsection*{Your turn}
\begin{itemize}
\raggedright
  \item \emph{Say it:} \emph{pośāk}
  \item \emph{Say it:} \emph{pośāk}, once more
  \item \emph{Say it:} \emph{mele meṁ bacce naī pośāk pahankar āe.} — The children came to the fair in new clothes.
  \item \emph{Recall:} say the Hindi for loose coins, then the Hindi for a size, then say \emph{pośāk} again
\end{itemize}

\begin{tcolorbox}[breakable,colback=teal!4,colframe=teal!35!black,title={Before you move on}]
What does \hi{पोशाक} mean? (\textquotedblleft{}A dress, an outfit\textquotedblright{}.) Say it once more.
\end{tcolorbox}
\section[\texorpdfstring{\hi{रुमाल} (rumāl)}{रुमाल (rumāl)}]{\hi{रुमाल} (rumāl) — a handkerchief}

\label{lesson:HI-C435-rumal}



\begin{quote}
Before the new one: say the Hindi for a size, then the Hindi for an outfit.
\end{quote}

\subsection*{You'll want to know: \hi{रुमाल}}
\textbf{\hi{रुमाल}} — \emph{rumāl} — \textquotedblleft{}a handkerchief\textquotedblright{}.

\begin{quote}
\hi{मैंने} \hi{बाज़ार} \hi{से} \hi{दो} \hi{रुमाल} \hi{ख़रीदे।}
\end{quote}

\begin{quote}
\emph{maiṁne bāzār se do rumāl kharīde.}
\end{quote}

\begin{quote}
I bought two handkerchiefs at the market.
\end{quote}

Masculine.

\begin{cousinweb}
\textbf{\hi{रुमाल}} (\emph{rumāl}) is usually traced to Persian \emph{rūmāl}, from \emph{rū}, \textquotedblleft{}face\textquotedblright{}, and a word for wiping: a cloth for the face.
\end{cousinweb}

\begin{grammarlens}[title={What you've built}]
One more word for getting around.
\end{grammarlens}

\subsection*{Your turn}
\begin{itemize}
\raggedright
  \item \emph{Say it:} \emph{rumāl}
  \item \emph{Say it:} \emph{rumāl}, once more
  \item \emph{Say it:} \emph{maiṁne bāzār se do rumāl kharīde.} — I bought two handkerchiefs at the market.
  \item \emph{Recall:} say the Hindi for a size, then the Hindi for an outfit, then say \emph{rumāl} again
\end{itemize}

\begin{tcolorbox}[breakable,colback=teal!4,colframe=teal!35!black,title={Before you move on}]
What does \hi{रुमाल} mean? (\textquotedblleft{}A handkerchief\textquotedblright{}.) Say it once more.
\end{tcolorbox}
\section[\texorpdfstring{\hi{सिमटना} (simaṭnā)}{सिमटना (simaṭnā)}]{\hi{सिमटना} (simaṭnā) — to huddle up; to draw together; to shrink}

\label{lesson:HI-C435-simatna}



\begin{quote}
Before the new one: say the Hindi for an outfit, then the Hindi for a handkerchief.
\end{quote}

\subsection*{You'll want to know: \hi{सिमटना}}
\textbf{\hi{सिमटना}} — \emph{simaṭnā} — \textquotedblleft{}to huddle up; to draw together; to shrink\textquotedblright{}.

\begin{quote}
\hi{ठंड} \hi{में} \hi{बच्चा} \hi{कंबल} \hi{में} \hi{सिमट} \hi{गया।}
\end{quote}

\begin{quote}
\emph{ṭhaṇḍ meṁ baccā kambal meṁ simaṭ gayā.}
\end{quote}

\begin{quote}
In the cold the child curled up in the blanket.
\end{quote}

Past (it …): \hi{वह} \hi{भीड़} \hi{एक} \hi{कोने} \hi{में} \hi{सिमट} \hi{गई}, \emph{vah bhīṛ ek kone meṁ simaṭ gaī}.

\begin{cousinweb}
\textbf{\hi{सिमटना}} (\emph{simaṭnā}) is a tadbhava of Sanskrit \emph{saṁveṣṭayati}, \textquotedblleft{}wraps up, winds round\textquotedblright{} (CDIAL 13026), which lists Hindi \emph{simaṭnā}, \textquotedblleft{}to be collected, shrink\textquotedblright{}; \hi{समेटना}, to gather up, is its partner.
\end{cousinweb}

\begin{grammarlens}[title={What you've built}]
One more word for getting around.
\end{grammarlens}

\subsection*{Your turn}
\begin{itemize}
\raggedright
  \item \emph{Say it:} \emph{simaṭnā}
  \item \emph{Say it:} \emph{simaṭnā}, once more
  \item \emph{Say it:} \emph{ṭhaṇḍ meṁ baccā kambal meṁ simaṭ gayā.} — In the cold the child curled up in the blanket.
  \item \emph{Recall:} say the Hindi for an outfit, then the Hindi for a handkerchief, then say \emph{simaṭnā} again
\end{itemize}

\begin{tcolorbox}[breakable,colback=teal!4,colframe=teal!35!black,title={Before you move on}]
What does \hi{सिमटना} mean? (\textquotedblleft{}To huddle up; to draw together; to shrink\textquotedblright{}.) Say it once more.
\end{tcolorbox}
